\section{Conclusiones, evaluaci\'on cr\'itica y pasos a seguir}
El proyecto implement\'o un flujo completo para recomendar portafolios en FinPUC, partiendo desde una base Markowitz, incorporando \textit{views} mediante Black-Litterman y cerrando con una simulaci\'on Monte Carlo P4 orientada a la decisi\'on cliente-empresa. Esta estructura permite evaluar no solo el desempe\~no financiero de cada cartera, sino tambi\'en su efecto en variables operativas relevantes: permanencia del cliente, aceptaci\'on de recomendaciones y utilidad esperada para la empresa.

La conclusi\'on principal es que Black-Litterman agrega valor cuando la estimaci\'on hist\'orica pura se vuelve inestable. La variante \texttt{momentum\_top20\_6m} es la m\'as consistente: con datos de pandemia alcanza Sharpe 2,38, superando a Markowitz (2,02) y al benchmark equiponderado (2,05). En la simulaci\'on P4 tambi\'en lidera el ranking, con Score 9.416,86, riqueza terminal esperada de USD 8.973,70, tasa de retiro de 3,1\% y utilidad empresa de USD 474,36. Esto la convierte en la mejor alternativa bajo el criterio combinado de valor para el cliente y rentabilidad para FinPUC.

El benchmark top-20 sigue siendo una referencia exigente, porque obtiene retorno bruto de 101,4\% sin depender de par\'ametros estimados. Sin embargo, no domina en retorno ajustado por riesgo. Markowitz, por su parte, es competitivo en el escenario sin pandemia, pero muestra mayor sensibilidad a la muestra de calibraci\'on. La validaci\'on robusta confirma este punto: el Sharpe futuro cae de 2,32 en h7\_f2 a 1,35 en h2\_f7, lo que muestra que el horizonte de entrenamiento incide materialmente en los resultados.

El principal punto d\'ebil est\'a en los supuestos estad\'isticos de la simulaci\'on. La normalidad se rechaza en los 601 activos, mientras que t-Student, Johnson SU y distribuciones relacionadas con colas pesadas ajustan mejor los retornos. Por eso, el Monte Carlo normal debe interpretarse como una aproximaci\'on inicial: es \'util para comparar t\'ecnicas bajo una misma regla, pero puede subestimar p\'erdidas extremas, CVaR y retiros en escenarios adversos. Tambi\'en quedan mejoras pendientes en costos de transacci\'on, calibraci\'on de la view macro de desempleo y definici\'on del gatillo de retiro del cliente.

\noindent\textbf{Fortalezas y puntos de mejora.}
\begin{center}
\small
\setlength{\tabcolsep}{4pt}
\renewcommand{\arraystretch}{1.05}
\begin{tabular}{@{}p{0.47\textwidth}p{0.47\textwidth}@{}}
\toprule
\textbf{Fortalezas} & \textbf{Puntos d\'ebiles y mejora} \\
\midrule
El dise\~no en tres capas permite medir el aporte de cada componente: optimizaci\'on, incorporaci\'on de views y simulaci\'on del comportamiento cliente-empresa. & El Monte Carlo usa retornos normales pese a la evidencia de colas pesadas; se propone probar t-Student o Johnson SU para capturar mejor riesgo extremo. \\
El benchmark top-20 es simple, reproducible y exigente, por lo que evita comparar los modelos contra una referencia d\'ebil. & Los costos de transacci\'on se consideran en P4, pero no dentro de la optimizaci\'on; se propone penalizar rotaci\'on o limitar $\|w_t-w_{t-1}\|_1$. \\
La comparaci\'on con y sin pandemia mantiene fijo el periodo fuera de muestra y a\'isla el efecto de la calibraci\'on hist\'orica. & La view de desempleo usa un supuesto fijo; debe reemplazarse por series macro observadas para evaluar capacidad predictiva real. \\
La validaci\'on con 70 splits, 18 ventanas m\'oviles, CSV de salida y semillas deterministas fortalece la trazabilidad del experimento. & $P_1$ mide p\'erdida contra $C_0$; usar drawdown desde el m\'aximo alcanzado ser\'ia m\'as cercano al comportamiento real del cliente. \\
\bottomrule
\end{tabular}
\end{center}

\noindent\textbf{Pasos a seguir.}
Para la etapa final, el foco debe estar en robustecer las partes que m\'as pueden cambiar la recomendaci\'on: distribuci\'on de retornos en Monte Carlo, sensibilidad de par\'ametros, costos de transacci\'on e integraci\'on operativa del recomendador. La Carta Gantt siguiente resume el plan.

\begin{center}
\scriptsize
\setlength{\tabcolsep}{3pt}
\renewcommand{\arraystretch}{1.05}
\begin{tabular}{@{}p{0.32\textwidth}ccccp{0.30\textwidth}@{}}
\toprule
\textbf{Actividad} & \textbf{12--19 May} & \textbf{20--31 May} & \textbf{1--15 Jun} & \textbf{16--26 Jun} & \textbf{Resultado esperado} \\
\midrule
Sensibilidad de $s$, $\tau$ y $\lambda$ & X & X &  &  & Identificar si el ranking cambia ante par\'ametros alternativos. \\
Monte Carlo con distribuciones no normales &  & X & X &  & Comparar riqueza, retiro y CVaR con colas m\'as realistas. \\
Costos de transacci\'on y rotaci\'on &  & X & X &  & Alinear optimizaci\'on y P4 mediante penalizaci\'on de turnover. \\
Integraci\'on del recomendador $<10$s &  &  & X & X & Prototipo ejecutable con respuesta operativa. \\
Informe, repositorio y presentaci\'on final &  &  & X & X & Documentaci\'on final y entrega del 26 de junio de 2026. \\
\bottomrule
\end{tabular}
\end{center}

En s\'intesis, la recomendaci\'on actual es usar Black-Litterman con view \texttt{momentum\_top20\_6m} como alternativa principal, manteniendo el benchmark como referencia y Markowitz como base metodol\'ogica. La decisi\'on final debe confirmarse con sensibilidad sobre distribuciones, costos y comportamiento del cliente, porque esos factores afectan directamente el Score P4 y la estabilidad del sistema recomendador.
